% 03-resultados.tex — Sección 3: resultados consolidados.
% Responde las preguntas 2 y 3 del revisor. Las cifras vienen de cifras/cifras.tex
% y el cuerpo de la tabla de tablas/resultados.tex: nada escrito a mano.

\section{Resultados consolidados}
\label{sec:resultados}

\subsection{Qué se consiguió medir}

La matriz son \cifraCeldas{} celdas (2 motores $\times$ 3 perfiles $\times$ 2
consultas $\times$ 3 concurrencias) con 3 repeticiones cada una:
\cifraFilasCsv{} filas, todas en \texttt{results/measurements.csv}, validadas
contra el esquema exacto de 14 columnas del manual. Ninguna combinación falta ni
está repetida, y se descartaron \cifraDescartadas{} repeticiones.

\begin{table}[htbp]
  \centering
  \caption{Reparto de las \cifraFilasCsv{} filas por \texttt{status}. La
           precedencia es \texttt{oom} > \texttt{error} > \texttt{spill} >
           \texttt{ok}, así que una fila \texttt{oom} puede tener además volcado a
           disco: por eso \cifraFilasConSpill{} filas tienen
           \texttt{disk\_spill = true} y solo \cifraFilasSpill{} salen con
           \texttt{status = spill}.}
  \label{tab:estados}
  \small
  \begin{tabular}{lrl}
    \toprule
    \texttt{status} & Filas & Qué significa \\
    \midrule
    \texttt{ok}    & \cifraFilasOk{}    & la celda se midió sin incidencias \\
    \texttt{spill} & \cifraFilasSpill{} & hubo volcado a disco durante la ventana \\
    \texttt{error} & \cifraFilasError{} & percentiles invalidados por censura (A4) \\
    \texttt{oom}   & \cifraFilasOom{}   & el kernel mató algún proceso del motor \\
    \bottomrule
  \end{tabular}
\end{table}

De las \cifraCeldas{} celdas, \cifraCeldasMedibles{} son \textbf{medibles} (sus
tres repeticiones salieron \texttt{ok} o \texttt{spill}). Las otras cuatro
contienen filas \texttt{error} u \texttt{oom}: en ellas el motor no completó
trabajo suficiente como para que un percentil signifique algo, y
\textbf{eso es el resultado}, no un fallo del instrumento.

\subsection{Tabla consolidada}

\begin{longtable}{llllrrrrrrll}
  \caption{Resultados por celda: media de las 3 repeticiones. Latencias en ms,
           throttling en segundos por ventana de 300~s, \texttt{max\_rss\_mb} en
           MiB.\label{tab:resultados}}\\
  \toprule
  Motor & Perfil & Consulta & c & qps & P50 & P95 & P99 & Throttl. & RSS & Spill & \texttt{status} \\
  \midrule
  \endfirsthead
  \toprule
  Motor & Perfil & Consulta & c & qps & P50 & P95 & P99 & Throttl. & RSS & Spill & \texttt{status} \\
  \midrule
  \endhead
  \bottomrule
  \endfoot
  % generado por scripts/build_report_data.py — no editar a mano
MongoDB & A & Q1 & 1 & 0.2911 & 3414 & 3743 & 4011 & 0.0 & 823.6 & no & ok \\
MongoDB & A & Q1 & 10 & 0.5444 & 18149 & 21824 & 23459 & 1180.8 & 872.9 & no & ok \\
MongoDB & A & Q1 & 50 & 0.6856 & 74756 & 83773 & 85571 & 1185.2 & 1040.0 & no & ok \\
MongoDB & A & Q2 & 1 & 0.7300 & 1352 & 1526 & 1704 & 0.0 & 831.9 & no & ok \\
MongoDB & A & Q2 & 10 & 1.9078 & 5110 & 6979 & 7779 & 1186.5 & 861.2 & no & ok \\
MongoDB & A & Q2 & 50 & 2.2711 & 21775 & 24751 & 26324 & 1196.1 & 961.9 & no & ok \\
MongoDB & B & Q1 & 1 & 0.2611 & 3818 & 4019 & 4167 & 12.8 & 812.2 & no & ok \\
MongoDB & B & Q1 & 10 & 0.3300 & 30597 & 35030 & 37341 & 1783.6 & 822.0 & no & ok \\
MongoDB & B & Q1 & 50 & 0.3333 & 138931 & 148438 & 151839 & 1788.0 & 879.4 & no & ok \\
MongoDB & B & Q2 & 1 & 0.1800 & 5536 & 5934 & 6130 & 17.4 & 829.8 & no & ok \\
MongoDB & B & Q2 & 10 & 0.4111 & 23692 & 26105 & 26166 & 1502.9 & 805.5 & no & ok \\
MongoDB & B & Q2 & 50 & 0.3733 & 146673 & 172727 & 179510 & 1773.4 & 908.8 & no & ok \\
MongoDB & C & Q1 & 1 & 0.1933 & 5505 & 5956 & 6082 & 158.2 & 587.7 & no & ok \\
MongoDB & C & Q1 & 10 & 0.1800 & 57199 & 63125 & 64473 & 2081.3 & 607.3 & no & ok \\
MongoDB & C & Q1 & 50 & 0.1356 & 283597 & 296988 & 298965 & 2085.2 & 634.7 & no & error\tnotaCensurada{} \\
MongoDB & C & Q2 & 1 & 0.0533 & 19329 & 20461 & 20683 & 183.4 & 611.6 & no & ok\tnotaBajaDensidad{} \\
MongoDB & C & Q2 & 10 & 0.1000 & 96905 & 113894 & 113895 & 1830.0 & 561.6 & no & ok \\
MongoDB & C & Q2 & 50 & 0.1667 & 225253 & 225372 & 228109 & 1976.7 & 672.3 & no & ok \\
PostgreSQL & A & Q1 & 1 & 0.1322 & 7481 & 8321 & 9030 & 0.3 & 103.8 & sí & spill \\
PostgreSQL & A & Q1 & 10 & 0.1444 & 74687 & 97188 & 98180 & 826.9 & 570.9 & sí & spill \\
PostgreSQL & A & Q1 & 50 & 0.0022 & 275776 & 283670 & 284371 & 552.6 & 1851.6 & sí & error\tnotaCensurada{} \\
PostgreSQL & A & Q2 & 1 & 1.8800 & 513 & 715 & 795 & 18.4 & 15.6 & no & ok \\
PostgreSQL & A & Q2 & 10 & 2.3511 & 3507 & 7477 & 8168 & 907.7 & 43.4 & no & ok \\
PostgreSQL & A & Q2 & 50 & 1.5789 & 35086 & 43086 & 50592 & 1093.6 & 127.3 & no & ok \\
PostgreSQL & B & Q1 & 1 & 0.1100 & 8923 & 10147 & 11198 & 146.0 & 105.0 & sí & spill \\
PostgreSQL & B & Q1 & 10 & 0.1000 & 114076 & 118444 & 119141 & 1535.2 & 571.0 & sí & spill \\
PostgreSQL & B & Q1 & 50 & 0.0133 & 208046 & 262901 & 270597 & 1216.0 & 1800.0 & sí & error/oom\tnotaCensurada{} \\
PostgreSQL & B & Q2 & 1 & 1.0889 & 885 & 1130 & 1311 & 550.0 & 16.3 & no & ok \\
PostgreSQL & B & Q2 & 10 & 1.1478 & 7330 & 15423 & 17939 & 1379.2 & 43.4 & no & ok \\
PostgreSQL & B & Q2 & 50 & 1.0556 & 51582 & 59091 & 62590 & 1634.3 & 127.3 & no & ok \\
PostgreSQL & C & Q1 & 1 & 0.0567 & 17555 & 18464 & 18777 & 489.3 & 104.4 & sí & spill\tnotaBajaDensidad{} \\
PostgreSQL & C & Q1 & 10 & 0.0333 & 263918 & 267051 & 267746 & 1829.8 & 564.7 & sí & spill\tnotaBajaDensidad{} \\
PostgreSQL & C & Q1 & 50 & 0.0000 & --- & --- & --- & 105.5 & 475.1 & no & oom\tnotaColapsada{} \\
PostgreSQL & C & Q2 & 1 & 0.5722 & 1730 & 2024 & 2108 & 870.7 & 16.7 & no & ok \\
PostgreSQL & C & Q2 & 10 & 0.6000 & 11936 & 28847 & 29888 & 1622.2 & 43.8 & no & ok \\
PostgreSQL & C & Q2 & 50 & 0.5611 & 100284 & 124632 & 134570 & 1890.9 & 127.0 & no & ok \\%

\end{longtable}

\leyendaNotas{}

\subsection{Throttling de CPU (pregunta 3 del revisor)}

Sí hubo throttling, y es la métrica que más fácil se malinterpreta. Dos
advertencias necesarias para leer la columna:

\begin{enumerate}
  \item \textbf{Puede superar la duración de la ventana.} El kernel suma el tiempo
        de throttling de \textbf{cada CPU} en la que el cgroup tenía hilos listos
        para correr. Con 4 CPUs y una cuota de 1, cada segundo de reloj puede
        acumular hasta 4 segundos de throttling. El máximo observado es de
        \textbf{2086~s en una ventana de 300~s}. No es un error de medida: cAdvisor
        reporta la misma suma, y las dos fuentes coinciden al 1~\%.
  \item \textbf{Una parte no la causa el benchmark.} Se midió el throttling de cada
        motor \textbf{en reposo}, sin consultas, por perfil: MongoDB acumula unos
        27~s por ventana de 300~s en el perfil~C y 3.6~s en el~B, mientras
        PostgreSQL se queda por debajo de 0.5~s en todos. Esa línea base está en
        \texttt{results/f6/f6\_3/baseline/} y aparece en la
        figura~\ref{fig:throttling-qps}.
\end{enumerate}

\subsection{Estabilidad de las mediciones}
\label{sec:estabilidad}

Antes de gastar las horas de la matriz se corrió una celda piloto completa y se
exigió que el coeficiente de variación del throughput entre las tres repeticiones
fuera menor del 15~\%; salió del 4.2~\% en MongoDB y del 3.8~\% en PostgreSQL.
Sobre la matriz entera, \textbf{la dispersión se mantuvo baja en casi todas las
celdas}, pero no en todas, y conviene decir cuáles:

\begin{itemize}
  \item \textbf{Dos celdas superan el 15~\%}, las dos de \texttt{Q2} con
        \mbox{c=50}: PostgreSQL en el perfil~A (20.9~\%, con una caída del 34~\%
        entre la primera y la tercera repetición) y MongoDB en el perfil~B
        (24.0~\%, con una caída del 33~\%).
  \item En la de \textbf{PostgreSQL} la causa es medible: las tres repeticiones
        gastaron \emph{el mismo} tiempo de CPU (1197~s) y el rendimiento del
        procesador cayó del 75.9~\% al 65.2~\% justo en la repetición que pierde el
        34~\%. El mismo CPU, más lento, hace menos trabajo.
  \item En la de \textbf{MongoDB} la causa \textbf{no está identificada}: el
        tiempo de CPU es idéntico (598.6~s) y el reloj está plano (83.4~\%,
        81.7~\%, 84.6~\%). La caché de Windows sobre el disco virtual, que
        \texttt{drop\_caches} no alcanza, es el sospechoso natural, pero no hay
        medición que lo demuestre y no se presenta como explicación.
  \item Se comprobó que \textbf{no hay un sesgo sistemático} de «la primera
        repetición es más rápida»: ocurre en 6 de las 30 celdas con las tres
        repeticiones comparables, cuando por azar serían unas 10.
\end{itemize}

La consecuencia práctica es que la celda de \texttt{Q2} con \mbox{c=50} en el
perfil~A es la \textbf{única} en la que MongoDB supera a PostgreSQL en
\texttt{Q2}, y es también una de las dos celdas con dispersión alta. Esa
comparación concreta no se usa para sostener ninguna conclusión
(sección~\ref{sec:interpretacion}).
